\section{Festigkeit: Woher die Kraft kommt und wie ein Tragwerk sie aufnimmt}\label{sec:festigkeit}

Die erste Festigkeitsrechnung der 1-GW-Anlage (Fusions-Surrogat mit 79 Proben, Hex-Modell und Elmer) ergab Spannungen im Spulengehäuse von
rund \SI{2000}{MPa}, viermal mehr als der Stahl bei \SI{4}{K} auf Dauer erträgt. Der Nutzer stellte dazu fest:
\emph{„Wand, Spulen und Festigkeit sind ein Thema, das gehört zusammen. Auch das Plasma gehört dazu. Spulenform
bestimmt Plasma, und die Spule muss auch halten, und das hat auch Einfluss auf die Temperatur.“} Später ergänzte der Nutzer
das Package (alles zwischen Plasma und Spule) und ein 3D-Modell als Ziel, und als die Zahlen kamen:
\emph{„Wir müssen uns auf die Festigkeit konzentrieren. 2000~MPa ist ein Killer, auch 500~MPa ist viel. Vielleicht
brauchen wir mehr Spulen, oder wir müssen verstehen, wie die Kraft entsteht.“} Dieser Abschnitt beschreibt, wie
daraus Schritt für Schritt ein Entwurf entstand, der die Grenzen einhält, unter welchen Bedingungen er gilt, und
warum er funktioniert. Er ist ein Zwischenstand: Die Entwurfsstudie (Schritt~5) rechnet noch.

\paragraph{Begriffe in diesem Abschnitt.}
\begin{description}\itemsep2pt
\item[Lorentzkraft] die Kraft, die ein Magnetfeld auf einen stromdurchflossenen Leiter ausübt. Je Meter Leiter ist
  sie $\vec f = I\,\vec t\times\vec B$ (Strom $I$, Richtung des Leiters $\vec t$, Feld $\vec B$). In unseren Spulen
  fließen \SI{22}{MA}; das ergibt bis zu \SI{140}{MN} je Meter, das Gewicht von \SI{14000}{t} auf jedem Meter.
\item[Eigenfeld] das Feld, das eine Spule an ihrem eigenen Leiter erzeugt. Es drückt die Spule auseinander wie einen
  aufgeblasenen Schlauch.
\item[Vergleichsspannung (von Mises)] eine Zahl, die den mehrachsigen Spannungszustand im Material auf einen Wert
  zusammenfasst, der mit der Festigkeit aus dem Zugversuch verglichen werden kann.
\item[99-\%-Wert (p99)] die Spannung, die in 99~\% des Stahlvolumens nicht überschritten wird. Einzelne Spitzen an
  Kanten des groben Rechennetzes werden so nicht überbewertet.
\item[Seilkraft] Ein gekrümmtes Seil kann eine Last, die quer auf es drückt, allein durch Zug tragen, wenn die Last
  zur Krümmung passt (wie ein Ring unter Innendruck). Ein gerades Stück kann das nicht: Es muss die Last durch
  \emph{Biegung} tragen, und Biegung erzeugt viel höhere Spannungen als Zug.
\item[Lager, Stützzylinder] die Stellen, an denen die Spulen festgehalten werden. Im Tokamak drückt jede Spule mit
  ihrem geraden inneren Schenkel gegen einen Stützzylinder in der Mitte der Anlage.
\item[Verbinder] Stahlblöcke zwischen benachbarten Spulengehäusen, die Kräfte von Spule zu Spule weitergeben.
\item[Dehnung] die relative Längenänderung. Das supraleitende Band (REBCO) verträgt dauerhaft nur einige Zehntel
  Prozent Zug, bevor es Schaden nimmt.
\item[Package] alle Schichten zwischen Plasma und Spule: Erstwand, Blanket (dort entsteht Tritium und Wärme),
  Abschirmung gegen Neutronen und Vakuumgefäß.
\item[Entwurfsstudie (DOE)] viele Entwürfe, deren Stellgrößen planmäßig über den ganzen Bereich verteilt werden
  (hier Latin Hypercube \cite{mckay1979}), damit man sieht, welche Größe was bewirkt.
\end{description}

\paragraph{Arbeitsteilung.} Der Nutzer gab die Richtung vor und entschied an jeder Gabelung. Die KI (Claude Opus)
plante, legte Erwartungen fest, schrieb die Kopplung und bewertete die Ergebnisse. Teilaufgaben gingen an kleinere
Modelle: Claude Sonnet recherchierte die Grenzwerte (jede Quelle per Crossref geprüft, Stichprobe von Opus
nachgeprüft) und erweiterte die Spulenoptimierung um ein Kraftziel; Claude Haiku erkundete, wie ein neuer
Spulensatz durch das Festigkeitsmodell läuft. Gerechnet haben simsopt \cite{simsopt} (Spulen) und Elmer
\cite{elmer} (Festigkeit).

\subsection*{Erwartung vorher}
Vor jeder Rechnung wurde im Ereignisprotokoll festgehalten, was erwartet wird und wie das Ergebnis zu deuten ist:
\begin{enumerate}\itemsep1pt
\item Ein Kraftziel in der Spulenoptimierung senkt die Spitzenkraft um etwa 30~\% bei doppeltem Feldfehler.
\item Kraftsenkung allein bringt nicht den nötigen Faktor~4; nötig ist ein Tragwerk, das die Kräfte geschlossen
  weiterleitet.
\item Die Kraft je Meter liegt bei 100--150~MN/m; die nötige Seilkraft schwankt entlang der Spule um mehr als den
  Faktor~3, weil die Krümmung stark wechselt. Biegung, nicht Zug, ist die Ursache der hohen Spannung.
\item Die Nettokraft auf jede Spule zeigt zur Mitte der Anlage.
\item Mehr Spulen senken die Kraft je Meter etwa umgekehrt proportional zur Spulenzahl; der Stahlbedarf bei gleicher
  Spannung bleibt etwa gleich.
\item Das Package ist mit dem jetzigen Abstand knapp.
\end{enumerate}

\subsection*{Schritt 1: Kraftziel in der Spulenoptimierung}
Die Spulenoptimierung (Stufe~2) bekam einen zusätzlichen Term, der hohe Kraft je Meter bestraft (simsopt
\texttt{LpCurveForce}). Ohne diesen Term entsteht exakt der bisherige Spulensatz (Abweichung des Feldfehlers
\SI{0,01}{\%}).
\begin{center}\small
\begin{tabular}{lccc}
\toprule
Gewicht Kraftziel & Feldfehler $B\cdot n$ (Mittel) & Spitzenkraft & mittlere Kraft \\
\midrule
0 (bisher) & 0,27~\% & 100~\% & 100~\% \\
0,01 & 0,30~\% & $-9$~\% & $-3$~\% \\
1 & 1,56~\% & $-25$~\% & $-6$~\% \\
\bottomrule
\end{tabular}
\end{center}
\textbf{Ergebnis:} Die Spitzen lassen sich kappen, die mittlere Kraft kaum. Sie ist durch Feldstärke und Größe der
Anlage festgelegt. Die Erwartung~1 war zu optimistisch, Erwartung~2 bestätigt: Die Spulenform allein löst das
Problem nicht. Das Kraftgewicht wird nur im Bereich 0 bis 0,03 weiterverwendet.

\subsection*{Schritt 2: Eine gekoppelte Rechenkette}
Damit Plasma, Spulen, Package, Festigkeit und Temperatur gemeinsam bewertet werden, wurde eine Kette gebaut
(\texttt{stellarator/kopplung.py}). Ein Entwurf durchläuft sie in 1--3~Minuten:
\begin{enumerate}\itemsep1pt
\item \textbf{Spulen} (simsopt): Spulenzahl, Abstand zum Plasma, Spulenlänge, Kraftziel $\to$ Feldfehler als Maß
  für die Plasmaqualität.
\item \textbf{Kalibrierung}: Strom so, dass auf der magnetischen Achse $B_0 = \SI{6,5}{T}$ herrscht. Wickelpaket
  quadratisch, Kantenlänge so, dass die Stromdichte wie bei 16 Spulen bleibt.
\item \textbf{Festigkeit} (Hex-Netz, Elmer): Gehäuse mit vier Wanddicken, Verguss, Verbinder, Lagerung; Last
  $\vec f = I\vec t\times\vec B$ aus allen Spulen.
\item \textbf{Package}: kleinster Abstand von der Gehäuse-Innenseite zum Plasmarand, rundum; daraus Abschirmdicke,
  Kernheizung und Neutronendosis im Wickelpaket.
\item \textbf{Temperatur}: Feld am Rand des Wickelpakets $\to$ Auslastung des Supraleiters und Temperaturreserve.
\end{enumerate}
Die Grenzen und Annahmen:
\begin{center}\small
\begin{tabular}{p{0.34\linewidth}p{0.18\linewidth}p{0.38\linewidth}}
\toprule
Größe & Grenze & Herkunft (Status) \\
\midrule
Gehäusestahl 316LN, p99 & \SI{500}{MPa} & 2/3 der Streckgrenze bei 4~K, ITER-Güte C3; beste Güte C1 erlaubt 667~MPa (ITER-Folien, ungeprüft) \\
Dehnung Wickelpaket & 0,3~\% & REBCO irreversibel ab 0,45--0,72~\% \cite{barth2015}; ARC hielt $<0{,}2$~\% \cite{sorbom2015} \\
Kernheizung Wickelpaket & \SI{2000}{W/m^3} & ARIES-CS (ungeprüft); streng \SI{50}{W/m^3} für EU-DEMO \cite{cufar2019} \\
Neutronendosis REBCO & $3\cdot10^{18}$~n/cm$^2$ & untere Schranke nach ARC \cite{sorbom2015} \\
Auslastung Supraleiter & 70~\% & Annahme; $J_c(B,T)$ nach \cite{senatore2024,pardo2026} \\
Feldfehler $B\cdot n$ & 0,3~\% & bisheriger Spulensatz als Maßstab \\
Package mindestens & \SI{1,0}{m} & ARIES-CS-Aufbau \cite{elguebaly2005} \\
Abschwächung Neutronen & e-fach je \SI{0,1}{m} & Annahme; Wolfram allein \SI{0,061}{m} \cite{windsor2021} \\
\bottomrule
\end{tabular}
\end{center}
Beim Bau der Kette fiel ein alter Fehler auf: Das bisherige Festigkeitsmodell hatte die Spulen um 3~\% größer
skaliert als Plasma und Alpharechnung. Die Kette nutzt jetzt überall denselben Maßstab. Der erste Testpunkt
(bisherige Spulen) bestätigte die Erwartungen bis auf eine: Das Feld am Leiter kam mit \SI{6,4}{T} statt etwa
\SI{15}{T} heraus. Ursache: Es war auf der Mittellinie gemessen, wo sich das Eigenfeld aufhebt. Am Rand des
Wickelpakets sind es \SI{18,7}{T}. Der Testpunkt zeigte: Package (\SI{1,77}{m} frei), Temperatur (Auslastung 34~\%)
und Feldgenauigkeit passen. Es scheitert allein an der Festigkeit (p99 \SI{2129}{MPa}) und an der Dehnung im
Wickelpaket (0,76~\%). Die Erwartung~6 (Package knapp) war falsch: ARIES-CS kommt mit \SI{1,3}{m} von Plasma bis
Spulenmitte aus \cite{elguebaly2005}, wir haben \SI{2,5}{m}.

\subsection*{Schritt 3: Woher die Kraft kommt}
Für jede Spule wurde die Kraft je Meter entlang des Leiters berechnet: das Feld der anderen Spulen nach Biot-Savart,
das Eigenfeld mit der regularisierten Formel für einen rechteckigen Querschnitt (simsopt). Zerlegt wird in einen
Anteil nach außen (die Spule weitet sich) und einen seitlichen Anteil.

\bild{kraft_analyse_16.png}{1.0}{Kraft entlang der vier Grundspulen (16 Spulen insgesamt, die übrigen sind
Spiegelbilder): Kraft je Meter, nötige Seilkraft, Krümmung.}
\lesehilfe{Jede Spalte ist eine Spule, waagrecht der Weg entlang der Spule (eine Spule ist etwa \SI{45}{m} lang).
Oben: Schwarz die gesamte Kraft je Meter, blau der Anteil, der die Spule aufweitet, rot der seitliche Anteil. Die
Kraft ist überall \SI{50}{MN/m} oder mehr und erreicht \SI{140}{MN/m} in der Mitte des Weges: Das ist der innere
Schenkel der Spule an der Innenseite des Rings, wo die Spulen eng beieinanderstehen und das Feld am stärksten ist.
Mitte: die Seilkraft, mit der die Spule diese Last als reiner Zug tragen könnte. Wo die Kurve fehlt, ist die Spule
fast gerade, dort ginge das nur mit unendlicher Seilkraft. Unten die Krümmung: Spitzen sind enge Bögen (Radius etwa
\SI{2,5}{m}), Täler sind fast gerade Stücke. Das Bild zeigt das Kernproblem: Die größte Kraft sitzt genau auf dem
geraden Stück.}

\textbf{Herleitung: Warum gerade Stücke so teuer sind.} Für ein Seil mit Seilkraft $T$, Krümmung $\kappa$ und
Querlast $f_N$ gilt im Gleichgewicht
\[
\frac{dT}{ds}\,\vec t + T\kappa\,\vec N + \vec f = 0 \quad\Rightarrow\quad T = \frac{f_N}{\kappa}.
\]
Mit $f_N \approx \SI{100}{MN/m}$ und $\kappa \approx \SI{0,25}{m^{-1}}$ braucht es $T \approx \SI{400}{MN}$. Bei rund
\SI{1,3}{m^2} Stahlquerschnitt im Gehäuse sind das schon etwa \SI{300}{MPa}. Das ist der Anteil, den man nicht vermeiden
kann, wenn Feld und Größe feststehen. Auf dem inneren Schenkel ist die Spule über etwa $l = \SI{7}{m}$ fast gerade
($\kappa < \SI{0,08}{m^{-1}}$), die Last dort ist $q \approx \SI{130}{MN/m}$. Als Balken auf zwei Stützen ergibt das
\[
M = \frac{q\,l^2}{8} \approx \SI{800}{MNm},\qquad
W = \frac{H^4 - h^4}{6H} \approx \SI{0,38}{m^3}\ (H = \SI{1,36}{m}\ \text{Gehäuse},\ h = \SI{0,755}{m}\ \text{Wickelpaket}),
\qquad \sigma_B = \frac{M}{W} \approx \SI{2100}{MPa}.
\]
Diese Handrechnung trifft den FEM-Wert von \SI{2129}{MPa} erstaunlich genau. Die hohen Spannungen kommen also aus
der Biegung des geraden inneren Schenkels.

\textbf{Wie Tokamaks das lösen.} Eine Tokamakspule hat die Form eines „D“: Der Bogen ist so geformt, dass er die Last
als reinen Zug trägt, und der gerade Schenkel innen liegt flächig an einem Stützzylinder oder an den Nachbarspulen an
(Verkeilung). Er wird also nicht gebogen, sondern gedrückt. Unser Modell dagegen hielt jede Spule nur an ihrer
äußersten Stelle fest. Die Nettokraft je Spule zeigt aber mit bis zu \SI{1,1}{GN} nach innen, zur Mitte der Anlage
(Erwartung~4 bestätigt). Die Spule hing damit wie ein Balken, der auf der falschen Seite gestützt ist.

\textbf{Mehr Spulen?} Dieselbe Analyse für 20 und 24 Spulen:
\begin{center}\small
\begin{tabular}{lcccc}
\toprule
Spulen & Strom je Spule & Kraft je Meter (Mittel) & Summe aller Spulen & Feldfehler $B\cdot n$ \\
\midrule
16 & \SI{22,1}{MA} & \SI{85}{MN/m} & \SI{58}{GN} & 0,27~\% \\
20 & \SI{17,7}{MA} & \SI{65}{MN/m} & \SI{57}{GN} & 0,074~\% \\
24 & \SI{14,7}{MA} & \SI{54}{MN/m} & \SI{58}{GN} & 0,028~\% \\
\bottomrule
\end{tabular}
\end{center}
Die Kraft je Meter sinkt, die Summe bleibt (Erwartung~5 bestätigt): Bei gleichem Stahlanteil folgt dieselbe
Spannung. Der Gewinn liegt woanders: Mit mehr Spulen wird das Feld zehnmal genauer. Diesen Spielraum kann man in
kürzere, glattere Spulen umtauschen.

\subsection*{Schritt 4: Tragwerksversuche mit denselben 16 Spulen}
Aus Schritt~3 folgte direkt, was zu ändern ist: dort abstützen, wohin die Kraft zeigt, und die Spulen zu einem
Verbund koppeln, der Querlasten von Spule zu Spule weitergibt. Das Festigkeitsmodell bekam dafür eine Lagerung an der
innersten Stelle jeder Spule (Fläche zur Anlagenmitte, wie ein Stützzylinder) und wahlweise innen und außen.

\bild{festigkeit_tragwerk.png}{1.0}{Gleiche 16 Spulen, anderes Tragwerk: Spannung, Verformung, Dehnung und Stahlmasse
für fünf Varianten (A--E), darunter die Spannung im Blick von oben für A und E.}
\lesehilfe{Oben vier Balkendiagramme für die fünf Varianten A bis E (Schlüssel in der Zeile darunter). Rote Balken
liegen über der Grenze (gestrichelt), blaue darunter, graue Größen haben keine feste Grenze. Von A nach E sinkt die
Spannung im Stahl von \SI{2129}{MPa} auf \SI{481}{MPa} und die Verformung von \SI{399}{mm} auf \SI{41}{mm}; dafür
steigt die Stahlmasse von \num{15600} auf \num{37000}~t. Unten links Variante~A von oben: rot sind Stellen über
\SI{1200}{MPa}, vor allem an den inneren Schenkeln um das Loch in der Mitte. Unten rechts Variante~E: fast alles blau
(unter \SI{400}{MPa}); die Spulen sind dicker und durch Verbinder zu einem Verbund geworden.}

\begin{center}\small
\begin{tabular}{llcccc}
\toprule
 & Variante & p99 Stahl & Verformung & Dehnung Wickelpaket & Stahl gesamt \\
\midrule
A & Lager außen (bisher), 6 Verbinder je Paar & \SI{2129}{MPa} & \SI{399}{mm} & 0,76~\% & \SI{15,6}{kt} \\
B & Lager innen & \SI{1454}{MPa} & \SI{273}{mm} & 0,55~\% & \SI{15,6}{kt} \\
C & innen und außen & \SI{1372}{MPa} & \SI{226}{mm} & 0,54~\% & \SI{15,6}{kt} \\
D & + 12 Verbinder je Paar, Stahl (200~GPa) & \SI{947}{MPa} & \SI{107}{mm} & 0,41~\% & \SI{23,1}{kt} \\
E & + dickste Wände des Bereichs & \SI{481}{MPa} & \SI{41}{mm} & 0,31~\% & \SI{37,0}{kt} \\
\bottomrule
\end{tabular}
\end{center}
\textbf{Warum das wirkt.} B: Die Nettokraft geht auf kurzem Weg in das Lager, statt die ganze Spule zu biegen
($-32$~\%, Erwartung 30--50~\% bestätigt). D: Mehr und steifere Verbinder geben die Querlast des geraden Schenkels an
die Nachbarn weiter; die freie Länge $l$ des Balkens wird kürzer, und $M \propto l^2$. E: Ein größeres
Widerstandsmoment $W$ senkt die Biegespannung weiter. Damit ist zum ersten Mal ein Entwurf unter \SI{500}{MPa}. Der
Preis ist die Masse. Das Problem hat sich verschoben: von „hält nicht“ zu „ist zu schwer“.

\subsection*{Schritt 5: Entwurfsstudie (Zwischenstand)}
Um den leichtesten zulässigen Entwurf zu finden, werden 200 Entwürfe gerechnet. Variiert werden elf Größen:
Spulenzahl (16, 20, 24), Abstand Spule--Plasma, Länge je Spule, Kraftziel, vier Wanddicken, Größe der Lagerflächen,
Zahl (6--14 je Paar) und Steifigkeit (20--200~GPa) der Verbinder. Fest sind Lager innen und außen. Erwartet wurden
10--25~\% zulässige Entwürfe, ein leichtester von 20--35~kt und die Dehnung als zweiter Engpass.

\bild{festigkeit_doe.png}{1.0}{Entwurfsstudie, Zwischenstand: Spannung über Stahlmasse (links) und über der Zahl der
Verbinder (rechts).}
\lesehilfe{Jeder Punkt ist ein vollständig gerechneter Entwurf. Links: waagrecht die Stahlmasse, senkrecht die
Spannung (logarithmisch, die gestrichelte Linie ist die Grenze von \SI{500}{MPa}). Die Farbe zeigt die Spulenzahl.
Schwarz umrandet sind Entwürfe, die \emph{alle} Grenzen einhalten: Festigkeit, Dehnung, Package, Kernheizung,
Supraleiter und Feldgenauigkeit. Rechts dieselben Entwürfe über der Zahl der Verbinder je Spulenpaar, die Farbe ist
die Steifigkeit der Verbinder. Masse allein entscheidet nicht: Es gibt schwere Entwürfe mit \SI{1000}{MPa} und
leichtere unter der Grenze.}

Bis zum Redaktionsschluss dieses Abschnitts waren 17 Entwürfe gerechnet, drei davon halten alle Grenzen ein:
\begin{center}\small
\begin{tabular}{lccc}
\toprule
 & Entwurf 8 & Entwurf 14 & Entwurf 17 \\
\midrule
Spulen & 24 & 20 & 24 \\
Länge je Spule & \SI{37}{m} & \SI{44}{m} & \SI{38}{m} \\
Abstand Spulenmitte--Plasma (Vorgabe) & \SI{2,13}{m} & \SI{2,35}{m} & \SI{2,22}{m} \\
Kraftgewicht & 0,030 & 0,025 & 0,019 \\
Verbinder je Paar, Steifigkeit & 11, \SI{142}{GPa} & 11, \SI{196}{GPa} & 7, \SI{29}{GPa} \\
Wände innen / seitlich & \SI{0,32}{m} / \SI{0,31}{m} & \SI{0,35}{m} / \SI{0,32}{m} & \SI{0,60}{m} / \SI{0,37}{m} \\
\midrule
p99 Stahl & \SI{305}{MPa} & \SI{377}{MPa} & \SI{389}{MPa} \\
Verformung & \SI{19}{mm} & \SI{18}{mm} & \SI{31}{mm} \\
Dehnung Wickelpaket & 0,17~\% & 0,21~\% & 0,20~\% \\
Stahl (Gehäuse + Verbinder) & \SI{25,9}{kt} & \SI{28,0}{kt} & \SI{22,6}{kt} \\
Feldfehler $B\cdot n$ & 0,22~\% & 0,22~\% & 0,17~\% \\
Freiraum Package & \SI{1,48}{m} & \SI{1,65}{m} & \SI{1,30}{m} \\
Kernheizung Wickelpaket & \SI{40}{W/m^3} & \SI{8}{W/m^3} & \SI{245}{W/m^3} \\
Feld am Leiter, Auslastung & \SI{15,7}{T}, 30~\% & \SI{16,9}{T}, 32~\% & \SI{16,0}{T}, 30~\% \\
\bottomrule
\end{tabular}
\end{center}

\textbf{Wie diese Entwürfe zustande kommen.} Alle drei haben 20 oder 24 Spulen, kürzere Spulen als bisher
(37--44~m statt 47~m) und ein Kraftziel nahe der Obergrenze. Sie nutzen also den Spielraum, den mehr Spulen bei der
Feldgenauigkeit schaffen, für kürzere und weniger stark gebogene Spulen. Zwei erreichen die Grenze über viele steife
Verbinder (geschlossener Verbund), Entwurf~17 stattdessen über eine dicke Innenwand (\SI{0,6}{m}) bei wenigen
weichen Verbindern. Die Dehnung im Wickelpaket bleibt bei allen deutlich unter 0,3~\%. Die Temperatur des
Supraleiters ist mit 30~\% Auslastung unkritisch. Entwurf~17 liegt näher am Plasma; seine Kernheizung erfüllt die
ARIES-CS-Grenze, nicht aber die strenge EU-DEMO-Grenze.

\subsection*{Vergleich mit der Erwartung}
\begin{center}\small
\begin{tabular}{p{0.45\linewidth}p{0.45\linewidth}}
\toprule
Erwartung & Ergebnis \\
\midrule
1. Kraftziel: $-30$~\% bei doppeltem Feldfehler & teurer: $-9$~\% bei $+12$~\%, $-25$~\% bei $\times5{,}7$ \\
2. Kraftsenkung allein reicht nicht & bestätigt \\
3. 100--150~MN/m, Seilkraft schwankt $>\times3$, Biegung ist die Ursache & bestätigt; Handrechnung trifft FEM \\
4. Nettokraft zur Anlagenmitte & bestätigt, bis \SI{1,1}{GN} je Spule \\
5. Mehr Spulen: gleiche Summe, gleiche Spannung & bestätigt; Gewinn ist die Feldgenauigkeit \\
6. Package knapp & widerlegt: 1,3--\SI{1,8}{m} frei, ARIES-CS braucht etwa \SI{1}{m} \\
Lager innen: $-30$ bis $-50$~\% & bestätigt: $-32$~\% \\
\bottomrule
\end{tabular}
\end{center}

\subsection*{Grenzen dieses Zwischenstands}
\begin{itemize}\itemsep1pt
\item \textbf{Lager starr.} Die Lagerflächen sind fest eingespannt, als wäre der Stützzylinder unendlich steif. Die
  zulässigen Entwürfe haben große Lagerflächen (Radius 2--3,5~m). Das ist optimistisch; der Stützzylinder selbst muss
  noch modelliert werden.
\item \textbf{Verbinder als volle Stahlblöcke.} Ihre Masse ist so groß wie die der Gehäuse. Ein Fachwerk oder eine
  Schale könnte leichter sein.
\item \textbf{Netz bei 24 Spulen gröber.} Damit der Speicher reicht, haben Entwürfe mit 20 und 24 Spulen weniger
  Stationen je Spule (77 bzw. 64 statt 96). Ein gröberes Netz kann Spitzen glätten. Die besten Entwürfe werden mit
  feinem Netz nachgerechnet, bevor sie als „hält“ gelten.
\item \textbf{Nur Primärlast.} Linear-elastisch, nur die Magnetkraft. Abkühlung, Eigengewicht, Ermüdung und
  Bruchmechanik (bei 4~K oft maßgebend) fehlen noch.
\item \textbf{Package und Temperatur grob.} Abschwächung der Neutronen und Bandanteil im Wickelpaket sind Annahmen;
  die Kernheizung reagiert empfindlich darauf.
\item \textbf{Plasma nur über den Feldfehler.} Die Alpha-Verluste werden erst für die besten Entwürfe voll gerechnet.
\end{itemize}

\subsection*{Nächste Schritte}
Entwurfsstudie zu Ende rechnen; Ersatzmodell trainieren und den leichtesten zulässigen Entwurf suchen; die besten
mit feinem Netz, elastischem Stützzylinder und Alpha-Bahnen nachrechnen; daraus das 3D-Modell in FreeCAD: Plasma,
Package-Schichten, Spulen mit Gehäuse und Tragwerk.
